
\section{Range States}
\label{Sect_range_states}

{\bf note} We will probably delete this section.

We will map the values of variables into their states.
Arrays will be mapped into states representing a multiset. 
We want to be able to handle in-place array conversions, where
one goes through an array with an index
and changes the properties of the elements.
In order to make this possible, it is not sufficient to 
keep track of single variables, because for a given index $ i, $ 
the values on positions before $ i $ may have different properties,
than the values on positions after $ i. $ 
We will therefore map pairs $ ( v, i ), $ where $ i $ is a potential
index of $ v $ separately. We will map $ (v,i) $ into a triple
of states, where the first element is the multiset of values before $ i, $ 
the second element represents the multiset of values on $ i, $
and the third element represents the multiset of values after $ i. $ 
Since we are possibly working on $ v[i], $ its state has to be stored
separately.

We will keep track of the
relative order between variables that are a potential index,
as well as $ 0 $ and the sizes of arrays.
The simplest ways to do this is by keeping a matrix of potential
differences.
If two indices are equal, their difference is zero.
If the first index is bigger than the second, the difference is positive.
\begin{definition}
   \label{Def_diff_set}
   A \emph{difference} has one of the following forms:
   $ -2^{*}, -1, 0, 1, 2^{*}. $
   A \emph{difference set} is a set of differences.
   We write $ \top $ for the full difference set
   $ \{ -2^{*}, -1, 0, 1, 2^{*} \}. $
\end{definition}
The intended meaning of $ -2^{*} $ is '-2 or less'.
The intended meaning of $ 2^{*} $ is '2 or more'.
Difference sets have two levels of uncertainty:
The first is in the semantics of $ -2^{*} $ and $ 2^{*}, $
which represent a set of possible difference values.
The second is the the fact that a set may contain more than one
element. This ugly asymmetry could be removed by replacing
$ -2^{*} $ by $ -2,-3,-4, \ldots, $ but we want the difference
sets to be finitely representable. We don't know a better
alternative to Definition~\ref{Def_diff_set}.

\begin{definition}
   We define two operations on difference sets:
   The first is \emph{negation}:
   If $ D $ is a difference set, then
   $ -D = \{ -d \mid d \in D \}, $
   using $ - $ on differences as defined in Figure~\ref{Fig_negation_oplus}.

   If $ D_1 $ and $ D_2 $ are difference sets, then
   $ D_1 + D_2 $ is defined as
   $ \bigcup \{ \, d_1 \oplus d_2 \mid
                   d_1 \in D_1 \mbox{ and } d_2 \in D_2 \, \}, $
   where $ \oplus $ is defined in Figure~\ref{Fig_negation_oplus}.
   Note that the entries in the table are sets, so that it is possible
   to take their union.
\end{definition}

\begin{figure}
   \caption{Unary $-$ and binary $ \oplus $}
   \label{Fig_negation_oplus}
   \[
      \begin{array}{l||c|c|c|c|c|}
         - & -2^{*} & -1 & 0 & 1 & 2^{*} \\
         \hline
         \hline
         & 2^{*} & 1 & 0 & -1 & -2^{*} \\
         \hline
      \end{array}
      \hspace*{1.0cm}
      \begin{array}{l||c|c|c|c|c|}
         \oplus & -2^{*} & -1 & 0 & 1 & 2^{*} \\
         \hline
         \hline
         -2^{*} & -2^{*} & -2^{*} & -2^{*} & -2^{*},-1 & \top \\
         \hline
         -1 & -2^{*} & -2^{*} & -1 & 0 & 1, 2^{*} \\
         \hline
         0 & -2^{*} & -1 & 0 & 1 & 2^{*} \\
         \hline
         1 & -2^{*}, -1 & 0 & 1 & 2^{*} & 2^{*} \\
         \hline
         2^{*} & \top & 1, 2^{*} & 2^{*} & 2^{*} & 2^{*} \\
         \hline
      \end{array}
   \]
   The entries in the table are always sets. Remember that $ \top $ denotes
   the complete set $ \{ -2^{*}, -1, 0, 1, 2^{*} \}. $
\end{figure}
\begin{definition}
   Let $ I $ be a finite set.
   A \emph{difference matrix} $ \Delta $ \emph{over} $ I $
   is a mapping of $ I \times I $ to
   difference sets.
   We write $ \Delta.{\bf at}(i_1,i_2) $ for
   $ \Delta(i_1,i_2). $
\end{definition}

Possible indices are variables that can have type $ {\bf u64}, $
as well as $ 0 $ and expressions of $ {\bf len}(s), $ where
$ s $ is a variable that holds an array term.
Note that the elements of $ I $ are symbolic: $ I $ consists
of expressions, not of their values.

We will now define a few operations on difference matrices.
We define them imperatively, instead of declaratively.
In other words, we define them as operations that modify
an existing difference matrix, instead of functions
that create a new one.
The reason for this choice is the following:
Below Definition~\ref{Def_range_state},
we will define the effect of a statement on a range state.
Due to overapproximation, and
the essential non-deterministic nature of high-level computation,
the effect of a statemenet is non-deterministic.
After having tried several
ways to specify the possible successor states, we concluded that
a description using non-deterministic computation is the simplest.
In an earlier version, we gave a recursive definition of the 
effect relation $ \vdash $ by means of Horn clauses,
but it turned out rather unreadable. Using imperative language
is somewhat natural, because in the end we are modeling an imperative
language.  

\begin{definition}
        \label{Def_operations_on_difference_matrix}
   Let $ \Delta $ be a difference matrix over an index set $ I. $
   \begin{itemize}
   \item
      $ \Delta. {\bf contains}(i) $ returns true if
      $ \Delta $ has a difference set for $ i. $ 
   \item
      $ \Delta. {\bf at}(i) $ returns the difference
      set for $ i. $ 
   \item
      If $ i \not \in I, $ then
      $ \Delta. {\bf extend}(i) $ modifies $ \Delta $ into
      a difference matrix $ \Delta' $ over $ I \cup \{ i \}, $ s.t.
      whenever both $ i_1, i_2 \not = i, $ then $ \Delta'. {\bf at}(i_1,i_2) =
      \Delta. {\bf at}( i_1,i_2). $
      If either of $ i_1, i_2 = i, $ then $ \Delta'.{\bf at}(i_1,i_2) = \top. $
   \item
      $ \Delta. {\bf assign}( i_1, i_2, d ) $ changes $ \Delta $
      into a difference matrix $ \Delta', $ s.t.
      $ \Delta'.{\bf at}(i_1,i_2) = d, $ and
      $ \Delta'.{\bf at}(i'_1,i'_2) = \Delta.{\bf at}(i'_1,i'_2), $
      whenever $ i'_1 \not = i_1 $ or $ i'_2 \not = i_2. $
   \item
      If $ i \in I, $ then
      $ \Delta. {\bf erase}(i) $ is obtained from $ \Delta $
      by making $ \Delta.{\bf at}(i,i_2) $ and $ \Delta.{\bf at}(i_1,i) $
      undefined, for all $ i_1,i_2 \in I. $ 
   \item
      $ \Delta. {\bf isconsistent}( ) $ returns true if
      there are no $ i_1, i_2 \in I, $ s.t.
      $ \Delta. {\bf at}(i_1,i_2) = \emptyset. $
   \end{itemize}
\end{definition}

\begin{definition}
   Let $ \Delta $ be a difference matrix. We define
   a procedure $ {\bf close}( \Delta ) $ that
   closes $ \Delta $ under reflexivity, transitivity and antisymmetry:
   $ {\bf close}( \Delta ) $
   is obtained by applying the following
   operations, until no further change is obtained:
   \begin{itemize}
   \item
      For every $ i, $ s.t. $ \Delta. {\bf contains}(i): $ \\
      \hspace*{1.0cm} $ 
         \begin{array}{l}
            \Delta. {\bf assign}(i,i, \{ 0 \} ); \hspace*{1.0cm}
            \Delta. {\bf assign}( 0, i, \,
               \Delta. {\bf at}(0,i) \cap \{ 0,1,2^{*} \} ); \\ 
            \Delta. {\bf assign}( i, {\bf len}(v),
               \Delta. {\bf at}( i, {\bf len}(v)) \cap \{ 0,1,2^{*} \} \, ) \\
         \end{array}
      $ 
   \item
      For every $ i_1,i_2, $ s.t. $ \Delta. {\bf contains}(i_1) $ and
      $ \Delta. {\bf contains}(i_2): $ \\
      \hspace*{1.0cm}
      $ 
           \Delta. {\bf assign}(i_1,i_2, \,
           \Delta. {\bf at}(i_1,i_2) \cap -\Delta. {\bf at}(i_2,i_1) \, ). 
      $
   \item
      For every $ i_1,i_2,i_3, $ s.t. $ \Delta. {\bf contains}(i_1), \
      \Delta. {\bf contains}(i_2), $ and $ \Delta. {\bf contains}(i_3): $ \\
      \hspace*{1.0cm} 
      $ \Delta.{\bf assign}(i_1,i_3), \
           \Delta. {\bf at}(i_1,i_3) \cap (
           \Delta. {\bf at}(i_1,i_2) + \Delta. {\bf at}(i_2,i_3)) \ ). $ 
   \end{itemize}
\end{definition}
We now define state sequences and a couple of operations
on them.
As for difference matrixes, we define the operations imperatively.
The operations on state sequences are not yet non-deterministic,
but they can fail.
\begin{definition}
   \label{Def_state_seq}
   Let $ {\cal A} = ( Q, \delta, N ) $ be a DMA.
   A \emph{state sequence} over $ {\cal A} $ is a
   (possibly empty) sequence of states of $ {\cal A}. $
   We define the following operations on state sequences:
   \begin{itemize}
   \item
      $ \overline{q}. {\bf push\_back}(q') $ replaces
      $ \overline{q} = (q_1, \ldots, q_n ) $ by
      $ (q_1, \ldots, q_n, q' ), $ and
      $ \overline{q}. {\bf push\_front}(q') $ replaces
      $ \overline{q} = (q_1, \ldots, q_n) $ by
      $ ( q', q_1, \ldots, q_n ). $
   \item
      If $ \overline{q} $ has form $ ( q_1, \ldots, q_n ) $ with $ n > 0, $
      then $ \overline{q}. {\bf pop\_back}( ) $ replaces
      $ \overline{q} $ by $ ( q_1, \ldots, q_{n-1} ), $
      and $ \overline{q}. {\bf pop\_front}( ) $ replaces
      $ \overline{q} $ by $ ( q_2, \ldots, q_n ). $
   \item
      If $ \overline{q} $ has form $ ( q_1, \ldots, q_n ) $ with $ n > 0, $
      then $ \overline{q}. {\bf back}( ) $ equals $ q_n, $
      and $ \overline{q}. {\bf front}( ) $ equals $ q_1. $
   \end{itemize}
\end{definition}
The purpose of a state sequence is to approximate the
data terms at positions $ v[i], v[i+1], \ldots. $ 
In fact, we currently do not use $ v[i+1], \ldots $ and looking at
$ v[i] $ is optional. Technically, we could use a 
single optional value. We defined Definition~\ref{Def_state_seq} a bit
more general, than needed. 

\begin{definition}
        \label{Def_operations_on_DMA_states}
   Let $ {\cal A} = ( Q, \delta, N ) $ be a DMA.
   \begin{itemize}
   \item
      Let $ q $ be a state of $ {\cal A}. $
      Operation $ q. {\bf insert}(q') $ replaces
      $ q $ by $ \delta( q;q' ). $
   \item
      Method $ q. {\bf extract}( ) $ non-deterministically
      does the following:
      Find $ q', q_1, $ s.t. $ \delta(q';q_1) = q. $
      Replace $ q $ by $ q' $ and return $ q_1. $
   \end{itemize}
\end{definition}
The second operation is one of the sources of non-determinism
in the effect of an operation, the other is uncertainty about the
relative order of indices.
\begin{definition}
   \label{Def_range_state}
   Let $ {\cal A} $
   be a an abstraction automaton.
   An \emph{range state} is a pair of form $ ( M, \Delta ), $
   s.t.
   \begin{itemize}
   \item
      $ M $ is a finite partial function, which
      \begin{itemize}
      \item
         maps single variables to $ Q. $
      \item
         maps pairs $ (v,i) $ consisting of two variables $ v,i $ to
         triples $ ( q_{1}, \overline{q}, q_{2} ). $
      \end{itemize}
      Both components can be partial.
   \item
      $ \Delta $ is a difference matrix over a set
      $ I $ consisting of variables, $ 0, $
      and expressions of form $ {\bf len}(v), $
      with $ v $ a variable.
   \end{itemize}
   We define the following methods of $ M: $
   \begin{itemize}
   \item
      $ M. {\bf contains}(v), \ M. {\bf contains}(v,i). $
      True if $ M $ contains a value for variable $ v, $
      or for the pair $ (v,i). $
   \item
      $ M. {\bf at}(v), \ M. {\bf at}(v,i). $ Returns the value
      for the variable $ v $ or the pair $ (v,i) $ in $ M. $ 
   \item
      $ M. {\bf assign}(v,d), \ M. {\bf assign}(v,i,d). $ Sets the value
      of variable $ v, $ or the pair $ (v,i) $ to $ d. $
   \item
      $ M. {\bf erase}(v), \ M. {\bf erase}(v,i). $ Remove the
      value for variable $ v, $ or the pair $ (v,d). $
   \end{itemize}
\end{definition}
The intended meaning of $ q_{1}, \overline{q}, q_{2} $
is that they represent the multisets of values
in $ v[ 0 \cdots i \rangle $, $ v[i] $, and $ v \langle i \cdots \rangle $, respectively.

For each program point, we will construct a set of
range states that are possible at that point,
using Definition~\ref{Def_semantics}.
As usual, it cannot be avoided that there will be abstract states
attached to a program point that are in fact unreachable.

\noindent
For simplicity, we will always assume that all assignments to
variables (not to a member) are initializations.
We explain the relation between assignment and
comparison. Since our computational model is non-deterministic,
assignment can be viewed as a special case of comparison,
where nothing is known about the assigned variable.
More precisely, the assignment $ i \leftarrow 0 $ can be viewed
as two statements: {\bf (1)} Initialize $ i $ with a random value.
{\bf (2)} Verify that $ i $ equals $ 0, $ failing if it is not.
Since computation is non-deterministic, $ i $ will be equal
to $ 0 $ in a succesful computation.
It is convenient to treat assignment as 
comparison as much as possible, and to deal with
known equalities in a postprocessing state. This has the advantage
that the common parts of assignment and comparison can be treated
by a single procedure.

We list possible statements, and their effect on range states.
Due to lack of space, the list is incomplete. 
\begin{itemize}
\item
   The statement $ i \leftarrow 0 $ has the following effect: \\
   \hspace*{1.0cm}
   $ 
      \Delta. {\bf extend}(i); \hspace*{0.5cm} 
      \Delta. {\bf assign}(i,0, \{ 0 \} ); \hspace*{0.5cm}
      {\bf close}( \Delta ); \hspace*{0.5cm}
      M. {\bf assign}( i, Q_{\cal A}(0) );
   $ 
\item
   A statement of form $ i \leftarrow {\bf len}(v) $
   has the following effect: \\
   \hspace*{1.0cm} 
   $ 
         \Delta. {\bf extend}(i); \hspace*{0.5cm}
         \Delta. {\bf assign}(i, {\bf len}(v), \{ 0 \} ); \hspace*{0.5cm}
         {\bf close}( \Delta );
   $
\item
   A statement of form $ i \leftarrow j $ 
   has the following effect: \\
   \hspace*{1.0cm} 
   $
         \Delta. {\bf extend}(i); \hspace*{0.5cm}
         \Delta. {\bf assign}(i, j, \{ 0 \} ); \hspace*{0.5cm}
         {\bf close}( \Delta ); 
   $
\item
   A statement of form $ w \leftarrow v[i] $
   has the following effect: \\
   \hspace*{1.0cm} 
   $ \begin{array}{l}
      {\rm let} \ ( q_1, \overline{q}, q_2 ) = M. {\bf at}(v,i) \\
      \mbox{if} \ \| \overline{q} \| = 0 \ \mbox{then} \\
      \hspace*{0.5cm} \overline{q}. {\bf push\_back}( q_2. {\bf extract}( )) \\
      M. {\bf assign}( w, \overline{q}. {\bf front}( )) \\
      M. {\bf assign}( v,i, ( q_1, \overline{q}, q_2 ) ) \\
      \end{array}
   $ \\ \\ 
   If the triple at $ M. {\bf at}(v,i) $ has a value for
   $ v[i], $ we can use it. Otherwise, we create it first. 
\item
   A statement of form $ j \leftarrow i + 1 $ has the following effect:
   \[
      \begin{array}{l}
      \Delta. {\bf assign}( j,i, \{ 1 \} ); \ \ {\bf close}( \Delta ) \\
      {\rm let} \ ( q_1, \overline{q}, q_2 ) = M. {\bf at}(v,i) \\
      \mbox{if} \ \| \overline{q} \| = 0 \ \mbox{then} \\
      \hspace*{0.5cm} \overline{q}. {\bf push\_back}( q_2. {\bf extract}( )) \\
      M. {\bf assign}( v,j, ( \delta( q_1; \overline{q}. {\bf front}()), q_2 )) \\ 
      \end{array}
   \] 
   We take $ M.{\bf at}(v,i), $ create the value at $ v[i] $ if it is not
   present, and transfer it to $ v[ 0 \cdots j \rangle. $ 
   Note that we are taking a slight shortcut.
   The problem is that $ i + 1 $ is not a flatterm.
   In reality, one must check for assignments of form
   $ j \leftarrow i_1 + i_2, $ where either
   $ \Delta.{\bf at}(i_1,0) = \{ 1 \}, $ or
   $ \Delta.{\bf at}(i_2,0) = \{ 1 \}. $ 
\item
   A statement of form $ v[i] \leftarrow w $ has the following effect:
   For every variable $ j $ that is a potential index of $ v, $ do: 
   \[ 
      \begin{array}{l} 
         {\rm require} \ \Delta(i, {\bf len}(v)) \cap \{ 0,1,2^{*} \} \not = \emptyset \\
         {\rm let} \ ( q_1, \overline{q}, q_2 ) = M. {\bf at}(v,j) \\
         {\rm nondeterministic} \\ 
         \hspace*{0.5cm} \bullet \ 
              {\rm require} \ \Delta.{\bf at}( i,j ) \cap \{ -1, -2^{*} \} \not = \emptyset \\
         \hspace*{0.8cm} q_1. {\bf extract}( ); \,
                         q_1. {\bf insert}( Q_{\cal A}( M. {\bf at}(w))); \\
         \hspace*{0.5cm} \bullet \ {\rm require} \
             0 \in \Delta.{\bf at}(i,j) \mbox{ and } 
                 \| \overline{q} \| = 1 \\
         \hspace*{0.8cm} {\rm let} \ \overline{q} = Q_{\cal A}( M. {\bf at}(w)) \\
         \hspace*{0.5cm} \bullet \  
          {\rm require} \ \Delta(i,j) \cap \{ 1,2^{*} \} \not = \emptyset 
                      \mbox{ or }
                          0 \in \Delta( i, j ) \mbox{ and } 
                          \| \overline{q} \| = 0 \\
          \hspace*{0.8cm} q_2. {\bf extract}( ); \,
                          q_2. {\bf insert}( Q_{\cal A}( M. {\bf at}(w))) \\
          M. {\bf assign}( v, i, \, ( q_1, \overline{q}, q_2 ) \, ) \\ 
       \end{array}
   \] 
   A require statement fails quietly if its condition is not met.
   The nondetermistic statement branches into three options 
   each of which can succeed or fail independently.
   The code may be surprising at first. First observe that $ i $ 
   will be included in the for loop. For $ i, $ the 
   middle option is entered if $ \overline{q} $ is non-empty,
   and $ \overline{q} $ will be replaced by 
   $ Q_{\cal A}( M. {\bf at}(w)). $ If $ \overline{q} $ is empty,
   the third option will be entered, and a random element from
   $ q_2 $ wil be replaced by $ Q_{\cal A}( M. {\bf at}(w)). $ 
   The assignment to $ v[i] $ may also affect 
   $ v[ 0, \ldots j \rangle $ or $ v[ j \ldots \rangle \rangle $ 
   for other variables $ j, $ for which the interval
   $ [ 0, \ldots j \rangle $ or $ [ j, \ldots, \rangle $ might
   contain $ i. $ 
\item
   A statement of form $ v \leftarrow w.f $ has the following effect:
   Let 
   $ i $ be the index of field $ f. $ 
   Let $ q = Q_{\cal A}(w). $  
   Non-deterministically pick a tuple $ ( q_1, \ldots, q_n ), $ for
   which $ \delta( q_1, \ldots, q_n ) = q. $
   $ M. {\bf assign}(v, q_i ). $ 
\item
   A statement of form $ v \leftarrow {\bf lt}(w_1,w_2) $
   has the following effect:
   \[
       \begin{array}{l}
          {\rm nondeterministic} \\
             \hspace*{0.5cm} \bullet \ \Delta. {\bf assign}( w_1, w_2,
               \Delta. {\bf at}( w_1, w_2 ) \cap \{ -2^{*}, -1 \} \, ); \ 
            {\bf close}( \Delta ) \\
            \hspace*{0.8cm} \mbox{require} \ \Delta. {\bf isconsistent}( )) \\
            \hspace*{0.8cm} M. {\bf assign}(v, Q_{\cal A}( {\bf t} )) \\
            \hspace*{0.5cm} \bullet \ \Delta. {\bf assign}( w_1, w_2,
               \Delta. {\bf at}( w_1, w_2 ) \cap \{ 0,1,2^{*} \} \, ); \ 
            {\bf close}( \Delta ) \\
            \hspace*{0.8cm} \mbox{require} \ \Delta. {\bf isconsistent}( )) \\
            \hspace*{0.8cm} M. {\bf assign}( v, Q_{\cal A}( {\bf f} )) \\
         \end{array}
      \]
      The other comparison operators are analogeous.
      Executing the code above results in either one or two successor states,
      dependent on whether assuming $ w_1 < w_2 $ or $ w_1 \geq w_2 $ is consistent
      with $ \Delta. $
\item
   A statement of form $ v \leftarrow F(w), $ with
   $ F $ a unary state condition has the following effect: \\
   \hspace*{1.0cm}
   $ \begin{array}{l}
         \mbox{if } F( Q_{\cal A}(w)) \\
         \hspace*{0.5cm} M. {\bf assign}(v, Q_{\cal A}({\bf t})) \\
         \mbox{else} \\
         \hspace*{0.5cm} M. {\bf assign}(v, Q_{\cal A}({\bf f})) \\
      \end{array}
   $ 
\end{itemize}
It remains to define a postprocessing procedure that must be called
after every operation on a range state.
Its purpose is to check for equalities inside the
difference matrix $ \Delta, $ and apply them. 
Similar to the description of the effect of a statement, we
describe the procedure in imperative fashion. 
\begin{definition}
   \label{Def_postproc}
   First define a function $ {\bf part}( q_1, \overline{q}, q_2 ) $ as
   follows: \\ 
   \hspace*{0.5cm}
   $ 
         \begin{array}{l}
            q = q_2; \ 
            \mbox{for} \ q' \in \overline{q} \mbox{ do} \\
            \hspace*{0.5cm}
               q = \delta( q'; q ) \\
            {\rm return} \ ( q_1, q ) 
         \end{array}
   $ \\ 
   The purpose of function $ {\bf part} $ is to partition
   the triple $ (q_1, \overline{q}, q_2) $ by pushing the middle component
   $ \overline{q} $ back into $ q_2. $ We use {\bf part} when
   comparing different $ M. {\bf at}(v,i), $ where we want
   to ignore the middle components.
   \begin{itemize}
   \item
      For every variable $ i, $ s.t.
      $ \Delta. {\bf at}(i,0) = \{ 0 \}, $ for every variable $ v, $
      s.t. $ M. {\bf contains}(v), $ do the following: 
      If $ M.{\bf contains}(v,i), $ then 
      require that $ {\bf part}( M. {\bf at}(v,i)) $ is equal to
      $ ( \delta( ), M. {\bf at}(v)). $ 
      Otherwise, $ M. {\bf assign}(v,i, ( \delta( ), ( ), M. {\bf at}(v)). $ 
   \item
      For every pair of variables $ i $ and $ v, $ s.t 
      $ \Delta. {\bf at}(i, {\bf end}(v)) = \{ 0 \} $ do the following:
      If $ M. {\bf contains}(v,i), $ then require that
      $ {\bf part}( M. {\bf at}(v,i)) $ is equal to
      $ ( M. {\bf at}(i), \delta()). $
      Otherwise,  
      $ M. {\bf assign}( v,i, ( M. {\bf at}(v), ( \, ), \delta( )). $ 
   \item
      For every pair of variables $ i, j, $ s.t.
      $ \Delta.{\bf at}(i,j) = \{ 0 \}, $ for every
      $ v, $ s.t. $ M. {\bf contains}(v), $ do the following:

      If $ M. {\bf contains}(v,i) $ holds, but not
         $ M. {\bf contains}(v,j), $ then 
      $ M. {\bf assign}(v,j, M. {\bf at}(v,i)). $ 
      Similarly, if
      $ M. {\bf contains}(v,j) $ holds, but
      $ M. {\bf contains}(v,i) $ does not hold, then
      $ M. {\bf assign}(v,i, M. {\bf at}(v,j)). $
      Otherwise, require that 
      $ {\bf part}( M. {\bf at}(v,i)) $ equals 
      $ {\bf part}( M. {\bf at}(v,j)). $ 
   \item
      For every variable $ v, $ s.t. $ M. {\bf contains}(v), $ 
      but $ v $ does not occur in $ {\bf live}_{\bf out}( \pi ), $
      call $ M. {\bf erase}(v). $
      For every pair of variables $ v,i, $ s.t.
      $ M. {\bf contains}(v,i), $ but
      either $ v $ or $ i $ does not occur in $ {\bf live}_{\bf out}(\pi), $
      call $ M. {\bf erase}(v,i). $
   \end{itemize}
\end{definition}
The description of the effects of statements that we gave is necessarily
incomplete, due to lack of space. In addition, there are many small 
variations possible. What we presented should be viewed as a toolbox
from which type checking procedures can be obtained.

We did not explain the type checking procedure itself. 
We give a short, intuitive overview. 
Initially, for every ambiguous statement, all candidate overloads
are inserted into the statement, joined by a {\bf resolve} statement.
We use Definition~\ref{Def_semantics} combined
with the range state transformations in Section~\ref{Sect_range_states}
to obtain a fixpoint containing combinations of possible states. 
After that, we cut out non-viable overloads 
from the flow graph and repeat the procedure until we reach
a fixed point where all overloads are viable. 
At this point, it must be the case that every
resolve statement is non-empty, and contain one candidate 
that is strictly more specific than all of the other candidates.
This can be checked with the procedure sketced below 
Definition~\ref{Def_acc_proj}.

